\documentclass{report}
\usepackage[T1]{fontenc} 
\usepackage[utf8]{inputenc} 
\usepackage{lmodern}
\usepackage [frenchb]{babel}
% Pour pouvoir utiliser 
\usepackage{ucs}

\usepackage{textcomp}
\usepackage{graphicx}
\usepackage{keystroke}
\usepackage{amssymb}
\usepackage{amsmath}
%\renewcommand{\thesection}{\arabic{section}} % numérotation des sections	
\usepackage[cc]{titlepic} %rajouter le logo dans la page de garde
\usepackage{url} % Pour avoir de belles url
\usepackage[a4paper]{geometry}
\usepackage[linktocpage]{hyperref}
% Pour pouvoir faire plusieurs colonnes
\usepackage {multicol}
%Pour faire plusieurs lignes
\usepackage{multirow}
\usepackage{slashbox}
% Pour mettre du code source
\usepackage {listings}
% Pour pouvoir passer en paysage
\usepackage{lscape}	

% Pour pouvoir faire plusieurs colonnes
\usepackage {multicol}

% POur crééer un index
\usepackage{makeidx}

%Pour taper des algorithme en pseudo code
%\usepackage{algorithm,algorithmic}

%insertion de pdf
\usepackage{pdfpages}


%Pour taper des algorithme en pseudo code
\usepackage{algorithm,algorithmic}
%Utilisation d'acronyme
\usepackage{acronym}

\usepackage{fancyhdr}
\pagestyle{fancy}
\usepackage{lastpage} %numérotation type k/n
\renewcommand{\footrulewidth}{0.5pt} % trait horizontal
%\renewcommand{\headrulewidth}{1pt} % suppresion du trait horizontal dans l'entête
\lfoot{Projet Services}
\cfoot{\thepage \hspace{0.15cm}sur  \pageref{LastPage} } % numérotation des pages
\rfoot{A.\textsc{M} \& R.\textsc{R}}
\fancyhead[L]{\hspace{-1.5cm}\includegraphics[scale=0.105]{img/logouniv.png}} %logo entête !


\fancypagestyle{plain}{% 1ères pages des chapitres:
  % \fancyhead{} % supprime l’entete...
  % \renewcommand{\headrulewidth}{0pt} % ...et le filet
}




\hypersetup{
  backref=true,
  %permet d'ajouter des liens dans...
  pagebackref=true,%...les bibliographies
  hyperindex=true, %ajoute des liens dans les index.
  colorlinks=true, %colorise les liens
  breaklinks=true, %permet le retour à la ligne dans les liens trop longs
  urlcolor= blue, %couleur des hyperliens
  citecolor=	cyan,
  bookmarks=true, %créé des signets pour Acrobat
  bookmarksopen=true,
  %si les signets Acrobat sont créés,
  %les afficher complètement.
  pdftitle={Rapport Projet Services}, %informations apparaissant dans
  pdfauthor={MARGUERITE Alain},
  %dans les informations du document
  pdfsubject={Doc}
  %sous Acrobat.
}

\makeindex
%%%% debut macro pour enlever le nom chapitre %%%%
\makeatletter
\def\@makechapterhead#1{%
  \vspace*{50\p@}%
          {\parindent \z@ \raggedright \normalfont
            \interlinepenalty\@M
            \ifnum \c@secnumdepth >\m@ne
            \Huge\bfseries \thechapter\quad
            \fi
            \Huge \bfseries #1\par\nobreak
            \vskip 40\p@
}}

\def\@makeschapterhead#1{%
  \vspace*{50\p@}%
          {\parindent \z@ \raggedright
            \normalfont
            \interlinepenalty\@M
            \Huge \bfseries  #1\par\nobreak
            \vskip 40\p@
}}
\makeatother
%%%% fin macro %%%%
%Couverture 
\widowpenalty=10000
\clubpenalty=10000
\hyphenation{permet-tant}

\title{ {\huge \textsc{Distributed System}} \\Master ALMA 2\up{eme} année \\\vspace{3cm} \emph{Causal Order Message Delivery : From Logical to Physical Time}\\{\small Encadrant : M.\textsc{Achour}}\vspace{3cm}}
\titlepic{ \includegraphics[scale=0.2]{img/logouniv}}
\author{A.\textsc{Marguerite}\\ R.\textsc{Rincé}\vspace{3cm}}
\date{Université de Nantes \\ 2 rue de la Houssinière, BP92208, F-44322 Nantes cedex 03, FRANCE}
%\vspace{20cm}


\newcommand{\cf}{\emph{cf}.}


\begin{document}

\maketitle
\renewcommand{\labelitemi}{$\bullet$} 

\clearpage

\tableofcontents
\clearpage

%\chapter*{Remerciements}
\phantomsection
\addcontentsline{toc}{section}{Remerciements}
 TODO TODO TODO TODO TODO TODO TODO TODO TODO TODO TODO TODO TODO TODO TODO TODO TODO TODO TODO TODO TODO TODO

 TODO TODO TODO TODO TODO TODO TODO TODO TODO TODO TODO TODO TODO TODO TODO TODO TODO TODO TODO TODO

\chapter{Introduction}\label{chap:Intro}
\section{Présentation du sujet}
L'objectif de ce projet/TP est de réaliser un protocole de transfert de messages entre deux machines avec des listes des canaux de communication \textit{First In First Out}. La démarche proposée par l'énnonce propose les  étapes  suivantes :
  
  
  \begin{enumerate}
  \item 
    Modéliser ce protocole avec l'outil \textsc{Roméo}.
  \item
    Construire le graphe des marquages.
  \item
    Spécifier en logique temporelle: \hfill \\
    le protocole peut toujours revenir à l’état initial.
  \item 
    Expliciter les propriétés vraies.
  \item
    Programmer une implémentation : \hfill \\
    Une machine est pilotée par l'utilisateur, l'autre par le programme.
  \item
    Tests  :\hfill \\
    rendre observable les différentes transitions.
  \end{enumerate}


\begin{algorithm}
\caption{\textbf{Recherche}(nœud $N$, rectangle $R$)}
/* \textit{Recherche l'ensemble des feuilles contenues dans un rectangle requête $d$-dimensionnel à partir d'un nœud $N$, les feuilles solutions seront stockées dans $\Lambda$} */
\label{algo:recherche}
\begin{algorithmic}[1]
\IF {$N$ n'est pas une feuille}
  \FORALL{éléments $n$ de $N$ dont le MBR intersecte $R$}
    \STATE \textbf{Recherche}($n$, $R$)
  \ENDFOR
\ELSE[$N$ est une feuille] %\COMMENT{$N$ est une feuille}
  \FORALL{éléments $n$ de $N$ dont le MBR intersecte $R$}
    \STATE ajouter $n$ à $\Lambda$
  \ENDFOR
\ENDIF
\RETURN $\Lambda$
\end{algorithmic}
\end{algorithm}

\begin{algorithm}
\caption{\textbf{Ajout}(élément $E$, nœud $NR$)}
/* \textit{Ajoute un nouvel élément $E$ dans le R-tree de nœud racine $NR$} */
\label{algo:ajout}
\begin{algorithmic}[1]
\STATE On parcourt l'arbre depuis la racine $NR$ jusqu'à la feuille appropriée. À chaque niveau, on sélectionne le nœud dont le MBR nécessite le plus petit élargissement pour contenir le MBR de $E$
\STATE En cas d'égalité, on choisit le nœud ayant le MBR de plus petite \og superficie \fg{}
\IF{la feuille atteinte $L$ n'est pas pleine}
  \STATE On insère $E$ dans $L$
  \STATE On met à jour tous les MBR de $L$ jusqu'à $NR$ pour qu'ils puissent contenir le MBR de E 
\ELSE[la feuille $L$ est pleine]
  \STATE On définie $\xi$ l'ensemble constitué de tous les éléments de $L$ et du nouvel élément~$E$\\
On choisit deux éléments $e_1$ et $e_2$ tel que la distance entre $e_1$ et $e_2$ soit la plus grande de toutes les autres paires de $\xi$\\
On crée deux nœuds $l_1$ et $l_2$ contenant respectivement $e_1$ et $e_2$
  \STATE On parcourt tous les éléments restants de $\xi$ et l'on assigne chacun d'eux à $l_1$ ou $l_2$, en fonction du MBR de ces nœuds qui permet le plus petit élargissement pour contenir le MBR de l'élément considéré
  \IF{égalité de l'élargissement}
    \STATE On assigne l'élément au nœud ayant le MBR de plus petite superficie
    \IF{égalité de superficie}
      \STATE On assigne l'élément au nœud ayant le moins d'élément
    \ENDIF
  \ENDIF
  \IF[On rappelle que m est le nombre minimal d'éléments que peut contenir un nœud]{il reste $\lambda$ éléments dans $\xi$ et qu'un nœud contient $m - \lambda$ éléments}
    \STATE On assigne tous les éléments restants dans $\xi$ à ce nœud
  \ENDIF
  \STATE On met à jour le MBR des nœuds se trouvant sur le chemin de $L$ à $NR$ pour que ceux ci couvrent le MBR de $l_1$ et $l_2$
  \STATE On effectue les \og splits \fg{} éventuels des nœuds supérieurs
\ENDIF 
\end{algorithmic}
\end{algorithm}


\chapter{Conclusion}\label{chap:COnc}
À ce stade, la modélisation de notre méta-modèle (\cf{} chapitre \ref{chap:M2}) nous a permis de générer un modèle pour notre application \textit{Clients-Serveur} (\cf{} chapitre \ref{chap:M1}). Par la suite il sera nécessaire de passer à l'étape 3 présentée au chapitre \ref{chap:Intro}. C'est à dire  d'implémenter le méta-modèle M2 en JAVA.

La création de l'application \textit{Clients-Serveur} repose sur l'héritage des classes du méta-modèle. Enfin l'instanciation concrète de ses classes achèvera l'implémentation de notre application.

% LocalWords:  méta-modèle Clients-Serveur l'instanciation
% LocalWords:  l'implémentation



\listoffigures
% Pour finir l'interligne de 1,5
%\bibliographystyle{alpha}
%\bibliography{biblio.bib}


\end{document}
